\section{Discussion}
\label{sec:discussion}

\paragraph{What generalizes.}
The specific numbers do not transfer beyond these two tasks; the design question
does, and our scale sweep sharpens what that question is. The naive reading of
the small-backbone result---``add a symbolic solver to your control group''---is
already refuted by our own 7B data, where the solver is 15 points behind and
contributes nothing. What survives both regimes is weaker and more portable: the
control group must be broad enough that \emph{some} arm in it is competitive
with RL under the conditions being reported, and it must be aggregated with a
$\max$ so that the answer cannot be improved by choosing the opponent.
Section~\ref{sec:statistics} quantifies the cost of neglecting this on our own
runs: the same 17 groups report $+15.07$ points against the warm start alone,
$+7.97$ against the mean of the control group, and $-4.44$ against its maximum.

\paragraph{What the recommendation is not.}
A reasonable objection is that ``aggregate with a $\max$ over a broad control
group'' is unfalsifiable as guidance, since a sufficiently determined critic can
always append one more arm until any result is negative. We agree that the
prescription needs a boundary, and we do not think this paper establishes where
it lies. Two properties distinguish our use of it from an open-ended ratchet.
The arm set was fixed by pre-registration before the confirmatory task ran, so
it could not grow in response to the result; and every arm receives the same
demonstration budget, so no arm can be strengthened by spending resources the
proposed method was denied. A protocol that fixes its arm set in advance and
equalizes budgets is a test. One that adds arms after seeing the outcome is not.
We make no claim to have found the right stopping rule for how broad a control
group must be.

\paragraph{Scope of the claim.}
Two tasks, three backbones spanning $14\times$, one RL algorithm (GRPO), one
task family (program-annotated financial QA), and deliberately scarce
demonstration budgets. Three narrower scope conditions deserve to be stated
plainly rather than inferred:

\emph{One algorithm, not a family.} Every RL number in this paper is GRPO with
one implementation. We have no evidence about PPO, RLOO, or the several GRPO
variants that correct its known estimator biases~\cite{drgrpo2025}. Where this
paper says ``RLVR'' as shorthand, the evidence underneath it is GRPO.

\emph{One parameterization.} Every arm---including every baseline---trains a
LoRA adapter at $r{=}16$, $\alpha{=}32$, learning rate $10^{-5}$. LoRA is
applied symmetrically, so it is not an unfair handicap on the RL arm
specifically, but it is a joint restriction on all arms and we cannot say
whether full-parameter training would change the ordering. It was not listed as
a threat in the conference version; it should have been.

\emph{The clean filter is our own, and it cannot manufacture the result.} An
item enters the study only if its gold program executes to its gold answer under
our verifier ($282/288$ of the CodeTAT-QA test set, $664/795$ of CodeFinQA). The
same verifier defines the training reward and grades the evaluation, which is
internally consistent but invites the charge that we selected the items on which
our conclusion holds. It cannot work that way. The filter is applied identically
to all five arms, so scoring every discarded item as wrong for everyone
multiplies each arm's accuracy by the same constant---$0.979$ for CodeTAT-QA,
$0.835$ for CodeFinQA---and a positive rescaling cannot change the sign of a
difference. We verified this directly: recomputing all 17 groups on the
unfiltered test sets leaves \textbf{17 of 17} signs unchanged, with $\Delta$
moving from $-4.44$ to $-4.00$ on average. No re-run was required; the objection
fails on arithmetic.

\emph{External calibration.} What the filter does affect is comparability of
absolute numbers, so we place ours against the published BizBench baselines,
which use the same $1\%$-relative-tolerance metric. On CodeTAT-QA our 7B arms
reach $79$--$85\%$, above the reported $75.0\%$ for Mistral-7B and $70.9\%$ for
CodeLlama-7B and in the band of CodeLlama-34B ($81.4\%$) and Mixtral-8x7B
($83.9\%$); our 0.5B GRPO reaches $63.12\%$, near Llama-2-13B ($65.1\%$). On
CodeFinQA our 1B arms reach up to $30.57\%$, comparable to MPT-30B ($31.0\%$)
and Llama-2-13B ($33.4\%$). These are not weak systems for their size. The
solver's $64.54\%$ sits between MPT-30B ($64.8\%$) and Llama-2-13B ($65.1\%$),
which is itself a quantification of this paper's point: 440 lines of rules match
a 13--30B model's few-shot performance on this task. One caveat about the
comparison itself: our splits are the public BizBench release
(\texttt{kensho/bizbench}), whose CodeFinQA and CodeTAT-QA test sets hold 795
and 288 items and whose training sets hold 4{,}668 and 2{,}856---all four
matching our copies exactly---whereas the BizBench paper's own table reports
test sets of 844 and 392. We cannot reconcile that difference from the outside,
so the published accuracies above should be read as the right order of magnitude
for each model scale rather than as strictly paired comparisons.

\paragraph{Threats to validity.}

\emph{The RL arm's budget is the largest threat, and it is not resolved.}
Section~\ref{sec:analysis} reports that no GRPO run completed one epoch over its
prompt pool ($0.10$--$0.79$), and that at 7B only 50 to 82 prompts per run
produced a non-zero group advantage. Both numbers point the same direction: our
GRPO arm is small by contemporary standards, and most visibly so at the scale
where this paper's generalization claim lives. The pre-registration fixed equal
wall clock, which at a fixed per-token cost buys an order of magnitude less data
at 7B than at 0.5B. We believe the 0.5B and 1B result is robust---there GRPO
learns vigorously, beating every learning arm by up to 35 points, and still
loses to the solver. We do not believe this study settles what a
well-resourced GRPO run does at 7B, and readers should not take it to.

\emph{Solver strength as a free parameter.} At 0.5B and 1B the solver sets the
bar. Section~\ref{subsec:solvercost} gives its development protocol, its dev
trajectory, its $7.46$-point dev-to-test penalty and its engineering cost, and
the logs are published. The 7B result also removes this objection structurally:
there the maximum is attained by continued SFT, an arm we did not design.
Section~\ref{sec:statistics} shows the verdict is negative in 8 of 17 groups
with the solver deleted from the control group entirely.

\emph{The operating points were chosen inside the scarce regime.} This is true
and it is by design: the pre-registered bands select the two smallest budgets
reaching stated warm-start accuracy windows, which places both operating points
well left of the $m \approx 1024$ crossing where learning arms overtake the
solver. The claim is therefore explicitly confined to the demonstration-scarce
regime. What we cannot offer is a confirmatory group past the crossing point: no
frozen-test evaluation was run at $m \geq 256$ in Config~1, so the statement
that the effect reverses there rests on dev dose curves rather than on the
decision rule.

\emph{Statistical power.} Under Holm correction across the family of 17, two
groups are significant. The per-group minimum detectable effect is a median of
$6.93$ points corrected, against measured effects with a median of $2.84$. The
quantitative claim rests on six independent operating points agreeing in sign
($p = 0.031$), not on per-group significance, and not on the count of 17.

\emph{Dev-selection bias is three observations, not a measurement.} We report
three instances where $\Delta$ is non-negative on dev and negative on the frozen
test. Each is real and the $\beta{=}0$ case is instructive, but three instances
identified after the fact do not establish the stated mechanism---that
policy-gradient trajectories fluctuate more on dev than fine-tuning trajectories
and so gain more from checkpoint selection. Measuring that would require
per-arm dev-checkpoint variance across all groups, which our stored artifacts
support only at 7B. We therefore state it as three observations consistent with
a mechanism rather than as a measurement of one, and a reader should weigh
contribution~(iii) accordingly: what is established is the practical point that
a frozen test set changes three verdicts, not the explanation for why.

\emph{Low absolute accuracy in Config~2.} The solver reaches only ${\sim}32\%$
there, so all arms operate in a low band. This is why the study spans a two-fold
difference in solver strength and a $14\times$ range in backbone; the sign is
unchanged across all of it.

\emph{One exclusion.} All 85 arm evaluations of the 17 reported groups are
present. The single excluded group is the collapsed run of
Section~\ref{sec:analysis}, whose $\Delta$ ($-19.73$) is more negative than any
we report, so excluding it weakens rather than strengthens the result. The
pre-registration did not specify an exclusion rule for training collapse; we
adopted one after observing it, and we report the result both ways.

\paragraph{What would change our mind, and what did.}
A null result should say what would refute it. We named three things. A GRPO run
at this scale with an order of magnitude more prompt coverage that beats
$\max\mathcal{B}$ at any operating point. A confirmatory group past the
$m \approx 1024$ crossing in which GRPO leads a control group including the
solver. Or a task family in the same program-annotated setting where no
deterministic arm is viable and GRPO still fails to beat continued SFT, which
would indicate our result is about the solver rather than about the control
group.

\textbf{The first one fired.} At ten times the prompt coverage, the
smallest-margin 7B cell reverses from $\Delta = -1.06$ to $+6.74$ and $+6.38$
across two training seeds, clearing our own Go threshold with McNemar
$p = 0.0013$ and $0.0029$ (Appendix~\ref{app:followup}). We report it in the
abstract, in Section~\ref{sec:scale}, in Section~\ref{sec:analysis} and here,
because a paper that specifies its own refutation conditions and then meets one
of them owes the reader the result in the same places it made the claim.

What survives is bounded rather than broken. The reversal is a budget effect,
not a defect effect: a matched-budget run that raises informative prompts from
82 to 120 lands at $+1.06$ with an interval of $[-2.13,+4.26]$, statistically
where the registered $-1.06$ already was. So under the pre-registered
protocol---equal wall clock for every arm---the finding holds at 7B; given ten
times that budget it does not. Config~3 therefore marks where this study's claim
ends, and the quantitative core is the 11 groups at 0.5B and 1B. The other two
tests we have not attempted.

\paragraph{On reporting negative results.}
The pre-registered protocol accepted both outcomes in advance, including
explicit language forbidding us from weakening the solver if GRPO won. We
mention this because a null result is only informative if the analysis could
have come out otherwise. In the same spirit, this extended version reports the
two findings that most damage our own case---two of seventeen groups significant
under our own correction, and an RL arm that never completed an epoch---in the
main text rather than in an appendix.

\section{Conclusion}

We pre-registered a decision rule that compares GRPO against the maximum of a
four-arm control group, one arm of which does not learn, and executed it on two
tasks and three backbones spanning $14\times$ in scale. All 17 frozen-test
effect sizes are negative, from $-1.06$ to $-16.42$ points. At 0.5B, the same
runs judged against a learning-only control group would have been reported as
substantial successes---GRPO beats every learning baseline there by up to 35
points.

The mechanism is not one strong baseline but a moving one. At 0.5B and 1B a
deterministic solver attains the maximum in all 11 groups, because it is flat in
the demonstration budget while every learning arm scales with it. At 7B the
solver falls 15 points behind and drops out of contention---exactly as the
``gains emerge with capability'' account predicts---and $\Delta$ remains negative
because continued SFT takes its place. Neither ablations (the solver is $100\%$
rule-derived) nor audits (zero constant-program shortcuts in 57{,}517 positive
samples) nor a hyperparameter sweep provides an escape; the sweep's one positive
$\Delta$ appears on dev and reverses on the frozen test.

This extended version adds the audit of our own instrument. The $\max$ in the
decision rule does not manufacture the result: a cross-fitted estimator that
selects the control arm on one half of the test set and scores it on the other
is negative in 17 of 17 groups. The sampling arm usually proposed as the missing
baseline is not deployable here, and even its oracle bound loses to the solver
at three of four operating points. Against that, only two groups survive Holm
correction, no GRPO run completed an epoch, and---most consequentially---the 7B
configuration reverses at ten times the budget, so it bounds the claim rather
than extending it. Within the registered protocol the negative sign is robust;
the per-group effect sizes are not; and the scale evidence is a budget boundary.

What generalizes is not the sign but the dependence. This study can now put
three numbers on how much a reported RLVR effect moves without any change to the
policy being evaluated: $19.5$ points for which control arm is nominated,
$7.8$ for an order of magnitude of compute, and $3.0$ for the inference engine
alone. An RLVR result reported without its comparison, its budget and its
decoding stack attached is not yet a result.

The practical recommendation is therefore not ``add a solver.'' It is: aggregate
your control group with a maximum rather than nominating one opponent, populate
it broadly enough that the aggregation can bind, fix the arm set before you see
the result, and reserve a test set that selection never touches. Which arm ends
up binding will depend on the regime---a solver in ours at small scale, plain
continued SFT at 7B---and that is precisely why the choice cannot be made in
advance by the party reporting the gain.
